\section{Introduction}

The last decade has seen a vast increase both in the diversity of applications to which machine learning is applied, and to the import of those applications. Machine learning is no longer just the engine behind ad placements and spam filters: it is now used to filter loan applicants, deploy police officers, and inform bail and parole decisions, amongst other things. The result has been a major concern for the potential for data driven methods to introduce and perpetuate discriminatory practices, and to otherwise be \emph{unfair}. And this concern has not been without reason: a steady stream of empirical findings has shown that data driven methods can unintentionally both encode existing human biases and introduce new ones (see e.g. \cite{swee13,bolukbasi2016man,CBN17,gendershades} for notable examples).

At the same time, the last two years have seen an unprecedented explosion in interest from the academic community in studying fairness and machine learning. ``Fairness and transparency'' transformed from a niche topic with a trickle of papers produced every year (at least since the work of \cite{PRT08}) to a major subfield of machine learning, complete with a dedicated archival conference (ACM FAT*). But despite the volume and velocity of published work, our understanding of the \emph{fundamental questions} related to fairness and machine learning remain in its infancy. What should fairness mean? What are the \emph{causes} that introduce unfairness in machine learning? How best should we modify our algorithms to avoid unfairness? And what are the corresponding tradeoffs with which we must grapple?

In March 2018, we convened a group of about fifty experts in Philadelphia, drawn from academia, industry, and government, to asses the state of our understanding of the fundamentals of the nascent \emph{science} of fairness in machine learning, and to identify the unanswered questions that seem the most pressing. By necessity, the aim of the workshop was not to comprehensively cover the vast growing field, much of which is empirical. Instead, the focus was on theoretical work aimed at providing a scientific foundation for understanding algorithmic bias. This document captures several of the key ideas and directions discussed.  